\scanpage{49}
\textcircled{\scriptsize 8} Soit \deleted{\unsure{$T$}} $\{\tau_i\}_{i\in I}\deleted{\unsure{\{\tau_1,\tau_2,\tau_3\}}}$ un triangle de triades biliées \inserted{(\unsure{noté|noté parfois|noté relation|relation non-liées} \unsure{Natif|Nat.|Notat.})}. \deleted{Alors}
\inserted{soit $I=\{i,j,k\}$,}
\deleted{Soit} Si $x_i\in\tau_i$, et si $x_j$ ($x_k$) est l'unique él. de
$\tau_j$ ($\tau_k$) non lié à $x_i$, alors $x_j$ et $x_k$ sont non liés.

En d'autres termes, si $J\subset\bigl(\prod_{i\in I}\tau_i,\deleted{\unsure{\tau_{i_2},\tau_{i_3}}}\bigr)$ est l'ens des $(x_i)_{i\in I}\deleted{\unsure{(x_j,x_k)}\text{ tels}}$ que $x_i$ non lié à $x_j$ pour $i\ne j$,
alors les projections $J\xrightarrow{\unsure{p_i|\mathrm{Pr}_i}}\tau_i$ sont bijectives (donc on a un syst. transitif d'isom entre les $\tau_i$, via \unsure{$J|t$}).

\deleted{\textcircled{\scriptsize 9} Soit \unsure{$S$ la} $X$ \inserted{$=\tau_1\cup\tau_2\cup\tau_3$} la réunion des trois triades d'un triangle $\{\tau_1,\tau_2,\tau_3\}$ de triades biliées $[\,\bar X=\bar\tau_1\cup\bar\tau_2\cup\bar\tau_3$ la réunion des trois triades du triangle opposé, alors $X,\bar X$ sont de cardinal $9$ et disjoints$]$.}

\textcircled{\scriptsize 9} Soit $X=\coprod_{i\in I}\tau_i\hookrightarrow S$, \inserted{(card $X=9$)} d'où $X\longrightarrow I$, dont les $3$ fibres sont \deleted{des} \inserted{les} \deleted{\unsure{$k$}} $3$ triades du triangle de triades donné:
\deleted{et} Soit $X\longrightarrow J$ l'application dont les restrictions aux $\tau_{\underline i}$ sont les $p_i^{-1}$, \deleted{d'où une application} \inserted{dont les fibres $(\tau'_j)_{j\in J}$ sont des $3$} autres triades \unsure{mutuellt|mutuellement} disjointes. Ceci posé

\textcircled{\scriptsize$\cdot$} $X$ est un bitriangle, \deleted{dont les} \inserted{dont} les six \deleted{\unsure{triades}} \inserted{triades}
(ou antitriangles) sont les $\tau_i$ \deleted{\unsure{$(i)$}} $(i\in I)$ $\tau'_j$ $(j\in J)$. forment les $2$ familles d'\deleted{anti}triades parallèles de $X$. Ainsi
\deleted{L'application $X\to I\times J$ est bijective.}
$I,J$ s'identifient resp. à l'une et l'autre familles de triades du bitriangle $X$, et l'application
\[
 X\longrightarrow I\times J
\]
est bijective.

\textcircled{\scriptsize 10} La donnée d'un triangle de triades biliées $(\tau_i)_{i\in I}$
\uline{équivaut} à celle d'un bitriangle $X$, muni
du \deleted{une} choix d'\uline{une} des deux familles \inserted{$I,J$} de
triades \unsure{mutuellmt|mutuellt|mutuellement} disjointes de $X$ (coorientation)
% archive sheet number: 48
